\documentclass[11pt]{article}
\usepackage[a4paper,margin=27mm]{geometry}
\usepackage{amsmath,amssymb,amsthm,mathtools,microtype,hyperref}
\usepackage[T1]{fontenc}
\usepackage{lmodern}
\hypersetup{hidelinks}
\newtheorem{theorem}{Theorem}
\newtheorem{lemma}[theorem]{Lemma}
\newtheorem{proposition}[theorem]{Proposition}
\theoremstyle{remark}
\newtheorem{remark}[theorem]{Remark}
\newcommand{\ddc}{dd^{\mathrm c}}
\newcommand{\dc}{d^{\mathrm c}}
\newcommand{\PSH}{\operatorname{PSH}}
\newcommand{\MA}{\operatorname{MA}}
\newcommand{\PP}{\mathbb P}
\title{Restricted masses and ordinary strict-current types\newline
on a nonclosed ruled surface}
\author{Research manuscript}
\date{October 7, 2026}
\begin{document}
\maketitle
\begin{center}
\small Working draft. Mathematical review and independent reconstruction pending.
\end{center}
\begin{abstract}
We propose restricted mass data for a concrete nonclosed background on a
ruled surface, with a logarithmic imposed divisor and an unbounded
nonanalytic residual. We optimize over every rough representative of the
ordinary singularity type and compute all intrinsic ideal trace orders in
a declared ambient embedding. The admissible type contains a strict current
exactly below an explicit divisorial threshold. Its two optimized masses
are preserved by pure pullback, whereas blowing up a point away from the
poles obstructs a strict current in the pulled-back ordinary type. A
one-sided strict witness still exists. The Gauduchon subtraction cost of
the residual background is computed separately. This is a proposed
restricted construction; the general Hermitian restricted-positivity
programme remains open.
\end{abstract}

\section{The construction}
Use $\ddc=(i/(2\pi))\partial\bar\partial$ and
$\dc=(i/(4\pi))(\bar\partial-\partial)$. Let $w$ be the probability
Fubini--Study form on $\PP^1$, and let $w_1,w_2$ denote its pullbacks to
$Y=\PP^1\times\PP^1$. For $p\in\PP^1$ let $G_p\leq0$ be the canonical
Green potential with $\ddc G_p=[p]-w$ and $\max G_p=0$. Write
\[
 D=\{p\}\times\PP^1,\qquad
 U=1-(1-G_p)^\alpha,\qquad
 r=bG_p+U,\qquad 0<\alpha<1,
\]
and pull these functions to $Y$. Fix $c>0$, $b\geq0$, and a smooth
strictly positive function $a$ on the first curve. Put
\[
 \theta=cw_1+aw_2.
\]
This background is strictly positive and is nonclosed when $a$ is
nonconstant. All functions use genuine canonical quasi-psh representatives.
When $b=0$, the function $bG_p$ is identically zero, including at $p$.
Ordinary equivalence $v\sim r$ means bounded difference, with the
inequalities retained at common poles. The function $r$ specifies a type;
its admissibility is a conclusion rather than an assumption.

For a nonempty admissible type define
\[
 R_-(\theta,[r])=\inf_{\substack{v\in\PSH(Y,\theta)\\v\sim r}}
                   \MA^{\rm ext}_\theta(v)(Y),\qquad
 R_+(\theta,[r])=\sup_{\substack{v\in\PSH(Y,\theta)\\v\sim r}}
                   \MA^{\rm ext}_\theta(v)(Y).
\]
The products are the finite-level non-pluripolar products with the smooth
background expansion. Every full-variable competitor is included.

\newpage
\begin{theorem}\label{thm:mass}
The specified ordinary type is nonempty precisely when $b<c$. In that
range it contains a genuine potential $v$ with
$\theta+\ddc v\geq\varepsilon(w_1+w_2)$ for some $\varepsilon>0$.
Every competitor has positive finite extended mass, and
\begin{align}
 \MA^{\rm ext}_\theta(v)(Y)
   &=\int_Y\theta^2-2b a(p)+2\int_Yv\,\ddc\theta,\label{eq:affine}\\
 R_-(\theta,[r])&=2(c-b)\min a,\qquad
 R_+(\theta,[r])=2(c-b)\max a.\label{eq:extrema}
\end{align}
Both extremal values are approached by actual representatives of that
ordinary type. No finite first-moment hypothesis is required.
\end{theorem}

This is an unbounded residual construction. It uses closed mass comparison
within its compact K\"ahler hypotheses \cite{BEGZ}; it does not assert
intersection invariance for the nonclosed form $\theta$.
The ordinary-type volume definitions and divisor and pure-pullback laws
under bounded mass are established background
\cite[Definition~3.1 and Proposition~3.3]{BGL}. Our purpose is their explicit
optimization in the declared nonanalytic family, including infinite first
moments. The companion manuscript \cite{Companion} records bounded-residual
trace formulas, divisorial normalization, residual Gauduchon calculations
and the local exceptional-curve obstruction. Their mechanisms retain their
existing credit; the obstruction below applies where our residual is bounded
near the chosen blowup center. No priority claim is made.

\section{Retaining the unbounded residual}
Let $\chi(t)=1-(1-t)^\alpha$ for $t\leq0$. Then
\[
 0<\chi'(t)=\alpha(1-t)^{\alpha-1}\leq\alpha,
 \qquad \chi''(t)=\alpha(1-\alpha)(1-t)^{\alpha-2}>0.
\]
The punctured-disk calculation gives
$\ddc U\geq-\alpha w$. The canonical removable-singularity theorem
\cite[I.5.23--5.24]{Demailly} extends this inequality across $p$.
The Lelong number is zero, since $U/\log|z-p|\to0$.

\begin{lemma}\label{lem:small}
There are genuine $U_L\sim U$ with $\sup U_L=0$,
$\ddc U_L\geq-\varepsilon_Lw$, $\varepsilon_L\to0$, and $U_L\to0$
in $L^1(\PP^1)$.
\end{lemma}
\begin{proof}
Put $\varepsilon_L=\chi'(-L)$ and
\[
 \chi_L(t)=
 \begin{cases}
  \chi(t)+C_L,&t<-L,\\
  \varepsilon_Lt,&-L\leq t\leq0,
 \end{cases}
 \qquad C_L=-\varepsilon_LL-\chi(-L).
\]
The values and derivatives match; $\chi_L$ is convex and increasing,
with $0\leq\chi_L'\leq\varepsilon_L$. The same curvature and canonical
extension argument applies to $U_L=\chi_L(G_p)$. Near $p$, $U_L-U=C_L$;
elsewhere the difference is bounded. Thus the mutual type and pole persist.
Writing $x=1+L$ gives
$C_L=(1-\alpha)x^\alpha+\alpha x^{\alpha-1}-1\geq0$.
Convexity also gives $\chi(t)\leq\varepsilon_Lt$ for $t\in[-L,0]$.
Hence $U\leq U_L\leq0$, and $U_L\to0$ pointwise off $p$.
Dominated convergence proves the $L^1$ assertion.
\end{proof}

\begin{lemma}\label{lem:remove}
Subtraction of the divisorial potential gives a canonical bijection
\[
 \{v\in\PSH(Y,\theta):v\sim r\}
 \longleftrightarrow
 \{q\in\PSH(Y,\beta):q\sim U\},\qquad
 \beta=(c-b)w_1+aw_2,
\]
and
\[
 \theta+\ddc(bG_p+q)=b[D]+\beta+\ddc q,
 \qquad \MA^{\rm ext}_\theta(bG_p+q)=\MA^{\rm ext}_\beta(q).
\]
\end{lemma}
\begin{proof}
Initially set $q=v-bG_p$ off $D$. In a chart write
$G_p=\log|z_1|^2-\sigma$, $\ddc\sigma=w_1$, and choose a quadratic
$\rho$ with $\ddc\rho\geq\theta$. The function
$v+\rho-b\log|z_1|^2=q+\rho-b\sigma$ is psh off $D$ and locally
bounded above, because $q\leq U+C$. Extend it canonically across $D$.
Its current retains the sharper lower bound $\ddc\rho-\theta$:
off $D$ that bound holds, and on $D$ its coefficient measures are positive
while the smooth correction has zero measure restriction. Subtracting
$\rho-b\sigma$ yields $\beta+\ddc q\geq0$.
The ordinary inequalities extend canonically; the upper bound forces the
same pole on $D$. Poincar\'e--Lelong gives the current identity. Conversely
that identity proves admissibility of $bG_p+q$. The two genuine functions
coincide canonically when they coincide almost everywhere.
Off $D$ their locally bounded residual products agree. Both non-pluripolar
products, including their smooth-background expansions, vanish on $D$.
This proves the measure identity.
\end{proof}

If $b<c$, choose $L$ with $\varepsilon_L<(c-b)/2$. Then
\[
 \theta+\ddc(bG_p+U_L)
 \geq b[D]+(c-b-\varepsilon_L)w_1+aw_2
 \geq\min((c-b)/2,\min a)(w_1+w_2).
\]
This is a strict witness in the mutual ordinary type. Conversely, average
the residual $q$ in Lemma~\ref{lem:remove} over the unitary group of the
second sphere. The background is invariant. Canonical Haar averaging gives
a genuine first-factor potential $q_0\sim U$ with
$(c-b)w+\ddc q_0\geq0$. Its total measure is $c-b$. Positivity excludes
$b>c$. If $b=c$, the measure is zero and the genuine harmonic $q_0$ is
constant on the compact curve, contradicting $q_0\sim U$.

\section{The full rough-class mass computation}
Suppose $k=c-b>0$ and put $\beta=kw_1+aw_2$. Choose $L_0$ with
$\varepsilon_{L_0}<k/2$, and use $u_0=U_{L_0}$ as an admissible residual
reference. Choose $C$ large enough that the closed K\"ahler form
$K=C(w_1+w_2)$ dominates $\beta$. The separated current
$S_0=K+\ddc u_0$ has no measure component on $D$: the first-factor
measure has zero atom at $p$. Consequently
\[
 \langle S_0\rangle=S_0,\qquad
 \langle S_0^2\rangle=2C(Cw_1+\ddc u_0)w_2,
 \qquad\int_Y\langle S_0^2\rangle=\int_YK^2.
\]
Now let $q\in\PSH(Y,\beta)$ be any competitor with $q\sim U$, and put
$S=K+\ddc q$. Both currents are closed and positive in the same K\"ahler
class. Their potentials have bounded difference and are locally bounded
off the closed complete pluripolar divisor $D$. BEGZ Theorem~1.16 applied
in both ordinary orders, for $p=1$ and $p=2$, gives
\[
 \{\langle S\rangle\}=\{K\},\qquad
 \int_Y\langle S^2\rangle=\int_YK^2.
\]
The positive polar current $S-\langle S\rangle=1_DS$ has zero trace
against $K$, and therefore vanishes. This proves $\langle S\rangle=S$.
All uses of closed cohomology concern $K$, rather than $\beta$.

Write $R=\beta-K$. Local bounded Bedford--Taylor expansion off $D$,
followed by finite trivial extension, gives
\[
 \MA^{\rm ext}_\beta(q)=\langle S^2\rangle+2R\wedge S+R^2.
\]
The terms are finite: the first has the closed mass just computed, the
second pairs a smooth form with a current of finite trace, and the last
is smooth. Distributional Stokes thus proves
\[
 \MA^{\rm ext}_\beta(q)(Y)=\int_Y\beta^2+2\int_Yq\,\ddc\beta.
\]
This pairing needs only $q\in L^1$, rather than a first moment against
its nonlinear measure. Haar averaging on the second factor preserves the
last expression and yields
\[
 \MA^{\rm ext}_\beta(q)(Y)=2\int_{\PP^1}a(kw+\ddc q_0).
\]
The curve measure is positive of total mass $k$, proving both bounds in
\eqref{eq:extrema} for every original competitor.

For exact extremizing sequences choose smooth strictly positive measures
$\zeta_j$ of mass $k$ converging to $k\delta_x$, and solve
$\ddc g_j=\zeta_j-kw$ on the sphere. Choose $L_j$ with
$\varepsilon_{L_j}$ below half the minimum density of $\zeta_j$ and
$|\int U_{L_j}\ddc a|<1/j$. Then $bG_p+g_j+U_{L_j}$ is a genuine
admissible representative of $r$, and its mass tends to $2k a(x)$.
Choosing $x$ at either extremum proves the formulas. Finally
$\int G_p\ddc a=a(p)-\int aw$ yields \eqref{eq:affine}.

The residual reference has infinite closed first moment when
$\alpha\geq1/2$. Indeed near $p$, writing $t=-G_p$,
$-U_{L_0}\asymp t^\alpha$ and
$\chi''(G_p)\asymp t^{\alpha-2}$. Angular integration of
$dG_p\wedge\dc G_p$ is comparable to $dt$. Its singular first moment
is comparable to $\int^\infty t^{2\alpha-2}\,dt$, finite exactly for
$\alpha<1/2$. The full mass proof above covers the divergent range too.

\section{Intrinsic ambient traces}
Embed the target as $Y_0=Y\times\{z_0\}$ in $Z=Y\times\PP^1$.
Use the strictly positive background
$\Theta=\operatorname{pr}_Y^*\theta+w_3$ and the genuine reference
$u_Z=\operatorname{pr}_Y^*(bG_p+U_L)$ chosen above.
Its generic Lelong number along $Y_0$ is zero. For the ambient divisor
$D_Z=D\times\PP^1$ and every real $s>0$,
\[
 \mathcal I(su_Z)=\mathcal O_Z(-\lfloor sb\rfloor D_Z),\qquad
 \mathcal I(su_Z)\mathcal O_{Y_0}
       =\mathcal O_{Y_0}(-\lfloor sb\rfloor D).
\]
To check the convention, factor a germ's largest transverse power
$z_1^j$. With $t=-\log|z_1|^2$, its radial integral has exponential
factor $\exp(-(j+1-sb)t+st^\alpha)$, up to bounded factors.
It converges exactly when $j+1>sb$, including divergence at equality.
A nonzero leading tangential coefficient gives necessity by Fubini;
local holomorphic bounds give sufficiency. No extra factor $2$ occurs
in the definition $|f|^2e^{-su_Z}$.

For every prime divisor $E$ on every smooth compact modification
$\pi:Y'\to Y_0$, orders of these restricted ambient principal ideals are
$\lfloor sb\rfloor\operatorname{ord}_E(\pi^*D)$. The intrinsic trace is
therefore
\[
 \tau_{(Y_0,u_Z)}(E)=b\operatorname{ord}_E(\pi^*D).
\]
This is the task's ambient-ideal definition. It asserts no commutation of
multiplier ideals with arbitrary pullback. The residual has trivial ideals
at every weight and is unbounded. At a generic point of $D$, an analytic
singularity of coefficient $b$ differs from $b\log|z_1|^2$ by a bounded
function; our residual diverges like $-t^\alpha$. Hence its ordinary type
is genuinely nonanalytic and is not determined by these ideal orders.

\section{Pure pullback and the strict-type obstruction}
Smooth gauges $(\theta,r,v)\mapsto(\theta+\ddc f,r-f,v-f)$ preserve
the actual currents, measures and both optimized masses. They preserve
multiplier ideals, since $e^{sf}$ is locally bounded above and below.

For a proper modification $\pi:Y'\to Y$, genuine
$\pi^*\theta$-psh functions descend canonically. Here is the precise
argument used for pure pullback. Add downstairs a local quadratic
$\rho$ with $\ddc\rho\geq\theta$. A potential upstairs plus
$\pi^*\rho$ is psh. Off the center it descends; properness gives a local
upper bound near the center, so canonical removability extends it.
Subtracting $\rho$ retains the sharper $\theta$ inequality: smooth
corrections have no measure restriction to the analytic center.
Pullback agrees almost everywhere, hence canonically. Ordinary bounded
inequalities descend in the same way. Off the exceptional set local
products agree by biholomorphism; non-pluripolar products vanish on the
excluded analytic sets. Thus
\[
 \pi_*\MA^{\rm ext}_{\pi^*\theta}(\pi^*v)=\MA^{\rm ext}_\theta(v).
\]
The full competitor bijection preserves both $R_-$ and $R_+$.
Exceptional curvature subtraction is a different background operation.
These are laws for the two ordinary-type masses; no invariance of
$V_\mu$ for arbitrary smooth positive upstairs probabilities is asserted.

\begin{proposition}\label{prop:blowup}
Blow up a point $y\notin D$. The pulled-back lower mass is positive and
there is a strict-current witness at most $\pi^*r$ plus a constant.
There is no strict-current witness ordinarily equivalent to $\pi^*r$.
\end{proposition}
\begin{proof}
Choose $q=U_L$ with $\beta+\ddc q\geq\delta(w_1+w_2)$.
For the exceptional divisor $E$, choose curvature $\beta_E$ of
$\mathcal O(E)$ restricting to negative Fubini--Study curvature on $E$.
For small $t>0$, $\pi^*(w_1+w_2)-t\beta_E$ is strictly positive.
For completeness, in finitely many blowup charts
$(x,y)\mapsto(x,xy)$ the pullback matrix $P$ has
$P_{x\bar x}\geq c_0$, $|P_{x\bar y}|\leq C|x|$, and Schur complement
at least $c_1|x|^2$. For $B=-\beta_E$ the vertical coefficient is
$B_{y\bar y}\geq b_0>0$ near $E$ and all entries are bounded.
The Schur complement of $P+tB$ is at least
\[
 c_1|x|^2+tb_0-C'(t|x|+t|x|^2+t^2).
\]
Fix the neighborhood small, then $t$ small, to bound the last terms by
$tb_0/2$. On the compact complement positivity persists under a small
bounded perturbation. The curvature metric extends smoothly from
$\mathcal O(E)|_E=\mathcal O_{\PP^1}(-1)$.
With $\ddc\ell_E=[E]-\beta_E$, the genuine potential
$\pi^*(bG_p+q)+\delta t\ell_E$ gives the strict lower bound
\[
 b[\pi^*D]+\delta(\pi^*(w_1+w_2)-t\beta_E)+\delta t[E].
\]
It is at most $\pi^*r$ plus a constant, but its exceptional pole changes
the ordinary type. The lower mass remains positive by pure pullback.

If $v'\sim\pi^*r$, it is locally bounded near $E$, since $y\notin D$.
Suppose $\ddc v'\geq\varepsilon h'-\pi^*\theta$ for some $h'>0$.
At each point of $E$ approximate the smooth lower bound from below by
a constant Hermitian form on a small adapted chart, with sufficiently
small scalar error. Its tangent coefficient is strictly positive, since
$\pi^*\theta|_E=0$. Subtracting a quadratic with this Hessian from $v'$
gives a genuine psh function. It restricts genuinely to $E$ because $v'$
is locally bounded. Adding back the quadratic proves that $v'|_E$ has
a strictly positive curve Hessian locally. A finite cover gives a positive
global lower bound against a probability form on $E$, contradicting
$\int_E\ddc(v'|_E)=0$. This proves the obstruction.
\end{proof}

\section{The separate Gauduchon cost}
Let $m>0$ be smooth on the first curve, $\int mw=1$, and put
$\mu=mw_1w_2$, $H=\int(m/a)w$. For all smooth positive tests $T$ with
$\ddc T=0$, set $D_\mu(T)=\int\sqrt{T^2/\mu}\,\mu$.
For $\beta=kw_1+aw_2$, $k=c-b>0$, the all-test volume is
\[
 V_\mu(\beta)=\left(\inf_T\frac{\int\beta T}{D_\mu(T)}\right)^2
      =\frac{2k}{H}.
\]
Indeed put $V=2k/H$ and $g=Vm/(2a)$, so $\int gw=k$. Solve
$\ddc f=(g-k)w$. Then $\Omega=gw_1+aw_2=\beta+\ddc f>0$ and
$\Omega^2=V\mu$. Stokes and the pointwise matrix arithmetic--geometric
mean inequality give $\int\beta T\geq\sqrt V D_\mu(T)$ for every test.
The closed positive test $T=(g/a)w_1+w_2$ attains it:
\[
 \int\beta T=2k,\qquad D_\mu(T)=\sqrt V H.
\]
This volume depends on $m$, whereas $R_\pm$ do not. It generally differs
from $R_+$; the original background has the separate volume $2c/H$.

Fix a real constant $d$ and a smooth real first-factor function $e$.
For $c_s=c-sd>b$, $a_s=a-se>0$, with fixed $b,m$, set
$k_s=c_s-b$, $H_s=\int(m/a_s)w$, $V_s=2k_s/H_s$ and use the preceding
minimizer $T_s$. Differentiation gives
\[
 H_s'=\int\frac{em}{a_s^2}w,\qquad
 \int(d w_1+e w_2)T_s=d+\frac{V_s}{2}H_s',\qquad
 \frac{d\sqrt{V_s}}{ds}
   =-\frac{\int(d w_1+e w_2)T_s}{D_\mu(T_s)}.
\]
Integration on every compact positive segment gives the actual residual
Gauduchon subtraction cost. It is not a derivative of $R_+$, and no
arbitrary smooth upstairs density is identified with a degenerate pullback.

\section{A concrete nonclosed example}
Let $\tau=(|z|^2-1)/(|z|^2+1)$ be the standard height on the first
sphere, with distribution $d\tau/2$ under $w$. Take
\[
 c=1,\qquad b=\frac14,\qquad \alpha=\frac34,\qquad
 a=2+\tau,\qquad m=1.
\]
For example $L=31$ gives
$\varepsilon_L=(3/4)32^{-1/4}<3/8$, so $bG_p+U_{31}$ is an actual
strict representative of the required nonanalytic type. Its residual has
infinite closed first moment. The full optimized masses and residual
Gauduchon volume are
\[
 R_-=\frac32,\qquad R_+=\frac92,\qquad
 H=\frac12\log3,\qquad V_\mu(\beta)=\frac3{\log3}.
\]
The original background instead has $V_\mu(\theta)=4/\log3$.
The ambient intrinsic orders are
$\tau(E)=\tfrac14\operatorname{ord}_E(\pi^*D)$ on every modification.
Blowing up any point outside $D$ preserves $R_-=3/2$ and $R_+=9/2$,
while Proposition~\ref{prop:blowup} excludes a strict current in the
pulled-back mutual ordinary type. The target backgrounds $\theta$, $\beta$,
$\Theta$ and their pure pullbacks are genuinely nonclosed; the auxiliary
$K$ is closed. The residual pole has not been replaced by a bounded or
analytic singularity.

\begin{remark}
Positive restricted lower mass supplies a mutual ordinary strict witness
downstairs within this declared family. It is preserved upstairs as a mass
invariant, while Proposition~\ref{prop:blowup} supplies only a one-sided
strict witness. The construction does not furnish a strict-type criterion
for arbitrary Hermitian backgrounds or solve the unrestricted S11 direction.
\end{remark}

\begin{thebibliography}{9}
\bibitem{BGL} S.~Boucksom, V.~Guedj and C.~H.~Lu,
\emph{Volumes of Bott--Chern Classes}, Peking Mathematical Journal (2025),
\href{https://doi.org/10.1007/s42543-025-00105-2}{doi:10.1007/s42543-025-00105-2},
especially Definition~3.1 and Proposition~3.3.
\href{https://d-nb.info/1374372846/34}{Published primary text}.
\bibitem{BEGZ} S.~Boucksom, P.~Eyssidieux, V.~Guedj and A.~Zeriahi,
\emph{Monge--Amp\`ere equations in big cohomology classes},
Acta Math. \textbf{205} (2010), 199--262, especially Definition~1.2,
Proposition~1.6 and Theorem~1.16.
\href{https://www.math.univ-toulouse.fr/~guedj/fichierspdf/BEGZ10.pdf}{Published primary text}.
\bibitem{Demailly} J.-P.~Demailly, \emph{Complex Analytic and Differential
Geometry}, Chapter~I, Theorems~5.23--5.24.
\href{https://www-fourier.univ-grenoble-alpes.fr/~demailly/manuscripts/agbook.pdf}{Author-hosted primary text}.
\bibitem{Companion} \emph{Non-K\"ahler pluripotential theory}, companion
manuscript: resolved trace masses, divisorial normalization and type
thresholds, divisorial Gauduchon volume and residual families.
\href{https://github.com/MingchenXia/non-kahler-pluripotential-theory/tree/1ff73204b833e197437c62dbe172815cb417bf7b/book/author}{Retained source version}.
\end{thebibliography}
\end{document}
